\documentclass[11pt,a4paper]{article}
\usepackage{fontspec,amsmath,amssymb,amsthm,microtype}
\setmainfont{lmroman10-regular.otf}[BoldFont=lmroman10-bold.otf,ItalicFont=lmroman10-italic.otf,BoldItalicFont=lmroman10-bolditalic.otf]
\usepackage{xcolor}
\usepackage[most]{tcolorbox}
\usepackage{mathrsfs}
\providecommand{\lt}{<}
\providecommand{\gt}{>}
\setmonofont{DejaVuSansMono.ttf}[BoldFont=DejaVuSansMono-Bold.ttf,ItalicFont=DejaVuSansMono-Oblique.ttf,BoldItalicFont=DejaVuSansMono-BoldOblique.ttf,Scale=0.86]
\newfontfamily\leanfallback{STIXTwoMath-Regular.otf}[Scale=0.86]
\definecolor{leanrule}{gray}{0.62}
\definecolor{leanback}{gray}{0.962}
\definecolor{leantitle}{gray}{0.22}
\definecolor{leanmuted}{gray}{0.38}
\newtcolorbox{leanbox}[1]{enhanced,breakable,colback=leanback,colframe=leanrule,
  boxrule=0.4pt,arc=0.6mm,left=2.4mm,right=2.4mm,top=1.8mm,bottom=1.8mm,
  before skip=0.8em,after skip=1em,fonttitle=\small\bfseries,coltitle=leantitle,
  colbacktitle=leanback,titlerule=0.3pt,title={#1}}
\newlength{\leancharwidth}
\newenvironment{leancode}{\par\smallskip\begingroup\ttfamily\footnotesize\raggedright
  \settowidth{\leancharwidth}{M}\setlength{\parskip}{0pt}\setlength{\parindent}{0pt}
  \frenchspacing}{\par\endgroup\smallskip}
\newcommand{\leanline}[2]{\hangindent=\dimexpr#1\leancharwidth+4\leancharwidth\relax
  \hangafter=1\noindent\hspace*{#1\leancharwidth}#2\par}
\newcommand{\leansource}[1]{\par{\footnotesize\color{leanmuted}#1}\par}
\newcommand{\leanlabel}[1]{\par\smallskip\noindent{\small\bfseries #1}\par}
\usepackage[margin=25mm]{geometry}
\usepackage[unicode,hidelinks]{hyperref}
\hypersetup{pdfauthor={LangevinLean Project}}
\setlength{\emergencystretch}{3em}
\setlength{\parskip}{0.4em}
\setlength{\parindent}{0pt}
\allowdisplaybreaks[2]
\raggedbottom
\title{\texorpdfstring{Strong limits of bounded increasing positive operators}{Strong limits of bounded increasing positive operators}}
\author{LangevinLean Project}
\date{}
\begin{document}\maketitle
Draft mathematical exposition; these notes may contain mistakes.


A sequence of positive operators with \(0\le A_n\le A_{n+1}\le cI\) has a bounded linear strong limit. This supplies an actual corrector operator from increasing finite sums, without choosing an eigenbasis.


\section{\texorpdfstring{The quadratic estimate behind convergence}{The quadratic estimate behind convergence}}


\textbf{Claim.} If \(0\le A\le cI\), then \(\|Ax\|^2\le c\langle Ax,x\rangle_{\mathbb R}\).



\textbf{Proof.} Positivity implies self-adjointness. The positive operators \(cI-A\) and \(A\) commute, so their product is positive. Its quadratic form is \(c\langle Ax,x\rangle-\langle A^2x,x\rangle\). Self-adjointness identifies the last term with \(\|Ax\|^2\). Rearrangement yields the bound.

\begin{leanbox}{Lean: \texttt{instance\_\allowbreak{}at\_\allowbreak{}20}, \texttt{positive\_\allowbreak{}operator\_\allowbreak{}square\_\allowbreak{}le}}
\leansource{\texttt{instance\_\allowbreak{}at\_\allowbreak{}20} \quad WassersteinMixing/\allowbreak{}MonotoneOperators.lean, lines 20–20}
\begin{leancode}
\leanline{0}{local\ instance\ :\ InnerProductSpace\ ℝ\ H\ :=\ InnerProductSpace.rclikeToReal\ ℂ\ H}
\end{leancode}
\leansource{\texttt{positive\_\allowbreak{}operator\_\allowbreak{}square\_\allowbreak{}le} \quad WassersteinMixing/\allowbreak{}MonotoneOperators.lean, lines 23–39}
\begin{leancode}
\leanline{0}{theorem\ positive\_operator\_square\_le\ (A\ :\ H\ →L[ℂ]\ H)\ (c\ :\ ℝ)}
\leanline{4}{(hA\ :\ 0\ ≤\ A)\ (hupper\ :\ A\ ≤\ c\ •\ (1\ :\ H\ →L[ℂ]\ H))\ (x\ :\ H)\ :}
\leanline{4}{‖A\ x‖\ \textasciicircum{}\ 2\ ≤\ c\ *\ inner\ ℝ\ (A\ x)\ x}
\end{leancode}
\leanlabel{How the Lean proof follows the argument}
The local scalar restriction gives the real part of the complex inner product. \texttt{positive\_\allowbreak{}operator\_\allowbreak{}square\_\allowbreak{}le} uses positivity of the commuting product and the proved symmetric pairing identity.
\end{leanbox}

\section{\texorpdfstring{Constructing the genuine bounded strong limit}{Constructing the genuine bounded strong limit}}


\textbf{Claim.} There is a bounded linear \(G\) with \(A_nx\to Gx\) for every \(x\), \(\|G\|\le c\), and \(0\le G\le cI\).



\textbf{Proof.} The increasing real numbers \(\langle A_nx,x\rangle\) are bounded by \(c\|x\|^2\), hence are Cauchy. For \(m\ge n\), apply the preceding estimate to \(A_m-A_n\); its scalar quadratic difference tending to zero forces \(A_nx\) to be Cauchy. Completeness gives the pointwise limit. Passing sums and scalar multiples through limits proves linearity, and the uniform norm bound proves continuity. Symmetry and the lower and upper quadratic inequalities pass to the limit.

\begin{leanbox}{Lean: \texttt{monotone\_\allowbreak{}operator\_\allowbreak{}cauchy}, \texttt{exists\_\allowbreak{}monotone\_\allowbreak{}operator\_\allowbreak{}limit}, \texttt{strong\_\allowbreak{}limit\_\allowbreak{}order}}
\leansource{\texttt{monotone\_\allowbreak{}operator\_\allowbreak{}cauchy} \quad WassersteinMixing/\allowbreak{}MonotoneOperators.lean, lines 43–89}
\begin{leancode}
\leanline{0}{theorem\ monotone\_operator\_cauchy\ (A\ :\ ℕ\ →\ H\ →L[ℂ]\ H)\ (c\ :\ ℝ)\ (hc\ :\ 0\ ≤\ c)}
\leanline{4}{(hmono\ :\ Monotone\ A)\ (hpos\ :\ ∀\ n,\ 0\ ≤\ A\ n)}
\leanline{4}{(hupper\ :\ ∀\ n,\ A\ n\ ≤\ c\ •\ (1\ :\ H\ →L[ℂ]\ H))\ (x\ :\ H)\ :}
\leanline{4}{CauchySeq\ (fun\ n\ ↦\ A\ n\ x)}
\end{leancode}
\leansource{\texttt{exists\_\allowbreak{}monotone\_\allowbreak{}operator\_\allowbreak{}limit} \quad WassersteinMixing/\allowbreak{}MonotoneOperators.lean, lines 93–119}
\begin{leancode}
\leanline{0}{theorem\ exists\_monotone\_operator\_limit\ (A\ :\ ℕ\ →\ H\ →L[ℂ]\ H)\ (c\ :\ ℝ)\ (hc\ :\ 0\ ≤\ c)}
\leanline{4}{(hmono\ :\ Monotone\ A)\ (hpos\ :\ ∀\ n,\ 0\ ≤\ A\ n)}
\leanline{4}{(hupper\ :\ ∀\ n,\ A\ n\ ≤\ c\ •\ (1\ :\ H\ →L[ℂ]\ H))\ :}
\leanline{4}{∃\ G\ :\ H\ →L[ℂ]\ H,\ ‖G‖\ ≤\ c\ ∧}
\leanline{6}{∀\ x,\ Tendsto\ (fun\ n\ ↦\ A\ n\ x)\ atTop\ ({\leanfallback 𝓝}\ (G\ x))}
\end{leancode}
\leansource{\texttt{strong\_\allowbreak{}limit\_\allowbreak{}order} \quad WassersteinMixing/\allowbreak{}MonotoneOperators.lean, lines 122–151}
\begin{leancode}
\leanline{0}{theorem\ strong\_limit\_order\ (A\ :\ ℕ\ →\ H\ →L[ℂ]\ H)\ (G\ :\ H\ →L[ℂ]\ H)\ (c\ :\ ℝ)}
\leanline{4}{(hpos\ :\ ∀\ n,\ 0\ ≤\ A\ n)\ (hupper\ :\ ∀\ n,\ A\ n\ ≤\ c\ •\ (1\ :\ H\ →L[ℂ]\ H))}
\leanline{4}{(hlim\ :\ ∀\ x,\ Tendsto\ (fun\ n\ ↦\ A\ n\ x)\ atTop\ ({\leanfallback 𝓝}\ (G\ x)))\ :}
\leanline{4}{0\ ≤\ G\ ∧\ G\ ≤\ c\ •\ (1\ :\ H\ →L[ℂ]\ H)}
\end{leancode}
\leanlabel{How the Lean proof follows the argument}
\texttt{monotone\_\allowbreak{}operator\_\allowbreak{}cauchy} uses monotone convergence of the scalar quadratic forms. \texttt{exists\_\allowbreak{}monotone\_\allowbreak{}operator\_\allowbreak{}limit} constructs the linear map from the actual pointwise limits. \texttt{strong\_\allowbreak{}limit\_\allowbreak{}order} transfers symmetry and both form inequalities through those limits.
\end{leanbox}

\section{\texorpdfstring{A commuting operator preserves the sharp energy factor}{A commuting operator preserves the sharp energy factor}}


\textbf{Claim.} If \(G\) is self-adjoint, \(\|G\|\le h\), \(B\ge0\), and \(GB=BG\), then

\[\langle Gx,BGx\rangle_{\mathbb R}\le h^2\langle x,Bx\rangle_{\mathbb R}.\]

\textbf{Proof.} The operator \(h^2I-G^2\) is positive because its quadratic form is \(h^2\|x\|^2-\|Gx\|^2\). It commutes with \(B\), so their product is positive. Its nonnegative quadratic form, together with symmetry and commutation, is exactly the displayed difference.

\begin{leanbox}{Lean: \texttt{commuting\_\allowbreak{}energy\_\allowbreak{}bound}}
\leansource{\texttt{commuting\_\allowbreak{}energy\_\allowbreak{}bound} \quad WassersteinMixing/\allowbreak{}MonotoneOperators.lean, lines 155–190}
\begin{leancode}
\leanline{0}{theorem\ commuting\_energy\_bound\ (G\ B\ :\ H\ →L[ℂ]\ H)\ (h\ :\ ℝ)\ (hh\ :\ 0\ ≤\ h)}
\leanline{4}{(hG\ :\ IsSelfAdjoint\ G)\ (hnorm\ :\ ‖G‖\ ≤\ h)\ (hB\ :\ 0\ ≤\ B)}
\leanline{4}{(hcomm\ :\ Commute\ G\ B)\ (x\ :\ H)\ :}
\leanline{4}{inner\ ℝ\ (G\ x)\ (B\ (G\ x))\ ≤\ h\ \textasciicircum{}\ 2\ *\ inner\ ℝ\ x\ (B\ x)}
\end{leancode}
\leanlabel{How the Lean proof follows the argument}
\texttt{commuting\_\allowbreak{}energy\_\allowbreak{}bound} constructs the positive squared-norm deficit, proves its commutation with the energy operator and expands the resulting positive product pairing.
\end{leanbox}

\end{document}
